\documentclass[11pt,a4paper]{article}
\usepackage[T1]{fontenc}
\usepackage[utf8]{inputenc}
\usepackage{lmodern,amsmath,amssymb}
\usepackage[margin=25mm]{geometry}
\usepackage{longtable,booktabs,array,calc}
\usepackage[hidelinks]{hyperref}
\hypersetup{pdftitle={Dobrushin's uniqueness condition for the Wilson action: a two-line strong-coupling gap at },pdfauthor={navstokgap research working notes}}
\newcommand{\tightlist}{\setlength{\itemsep}{0pt}\setlength{\parskip}{0pt}}
\setlength{\emergencystretch}{3em}
\setcounter{secnumdepth}{0}
\begin{document}
\hypertarget{dobrushins-uniqueness-condition-for-the-wilson-action-a-two-line-strong-coupling-gap-at-g2444-and-the-block-criterion-that-is-the-finite-verification}{%
\section{\texorpdfstring{Dobrushin's uniqueness condition for the Wilson
action: a two-line strong-coupling gap at \(g^2>444\), and the block
criterion that is the finite
verification}{Dobrushin's uniqueness condition for the Wilson action: a two-line strong-coupling gap at g\^{}2\textgreater444, and the block criterion that is the finite verification}}\label{dobrushins-uniqueness-condition-for-the-wilson-action-a-two-line-strong-coupling-gap-at-g2444-and-the-block-criterion-that-is-the-finite-verification}}

The finite-volume criteria of
\href{intermediate-region-finite-verification.md}{the
finite-verification note} have a single-site ancestor, Dobrushin's
uniqueness condition (Theory Probab. Appl. 13 (1968) 197), and for the
Wilson action it can be checked by hand. The conditional distribution of
one link given all the others is
\[d\mu_\ell(U\,|\,\omega)\ \propto\ \exp\Big[\frac{\beta_W}{N}\operatorname{Re}\operatorname{tr}\big(U\,M_\ell(\omega)\big)\Big]d\mu(U),
\qquad M_\ell=\sum_{p\ni\ell}S_p ,\] the sum of the six staples of
\(\ell\) in four dimensions. Changing one neighbouring link changes one
staple, and the resulting total-variation distance between the two
conditional distributions is at most \(e^{4\beta_W}-1\), independently
of \(N\). Each link has \(18\) neighbours that share a plaquette with
it, so Dobrushin's coefficient is
\[\alpha=\sup_\ell\sum_{\ell'}\rho_{\ell\ell'}\ \le\ 18\big(e^{4\beta_W}-1\big),\]
and the condition \(\alpha<1\) holds for
\(\beta_W<\tfrac14\log\tfrac{19}{18}=0.01352\), that is
\[g^2\ >\ 148\,N\ =\ 444\quad(SU(3)).\] Dobrushin's theorem gives a
unique Gibbs state, Künsch's theorem (Commun. Math. Phys. 84 (1982) 207)
gives exponential decay of all truncated correlations at rate
\(\log(1/\alpha)\) per lattice unit, and link reflection positivity
turns the rate into a gap:
\[\Delta_W\ \ge\ \frac{\hbar c}{a}\log\frac{1}{18(e^{24/g^2}-1)}\ \simeq\ \frac{\hbar c}{a}\log\frac{g^2}{432}\qquad(g^2>444),\]
uniformly in the volume. This is a better threshold than the polymer
expansion's \(1056\) and a smaller rate than its \(4\log(g^2/1056)\).
The block version of the same condition, the Dobrushin--Shlosman
constructive criterion (in \emph{Statistical Physics and Dynamical
Systems}, Birkhäuser 1985, 347), replaces the single link by a box \(V\)
of links and is expected to hold deeper into the intermediate region as
\(V\) grows; it is the quantity a finite verification would compute, and
Section 4 states it as a specification. All references at metadata
level; constants explicit; nothing promoted.

\hypertarget{the-one-link-conditional-distribution}{%
\subsection{1. The one-link conditional
distribution}\label{the-one-link-conditional-distribution}}

Fix all links except \(\ell\). The plaquettes containing \(\ell\) are
six in four dimensions, and each contributes
\(\operatorname{Re}\operatorname{tr}(U_\ell S_p)\) with \(S_p\in SU(N)\)
the product of the other three links of \(p\) in the appropriate order.
Hence the conditional density is \(e^{f_\omega(U)}\) with
\(f_\omega(U)=(\beta_W/N)\operatorname{Re}\operatorname{tr}(UM_\ell)\),
\(M_\ell=\sum_{p\ni\ell}S_p\), normalized by
\(Z(\omega)=\int e^{f_\omega}d\mu\). Gauge invariance of the interaction
is visible in the fact that a gauge transformation at an endpoint of
\(\ell\) rotates \(M_\ell\) and the Haar measure absorbs it; nothing
more is needed.

\hypertarget{the-dobrushin-coefficient}{%
\subsection{2. The Dobrushin
coefficient}\label{the-dobrushin-coefficient}}

\textbf{Lemma 1.} For probability densities \(p\propto e^f\),
\(q\propto e^{f'}\) with respect to the same measure,
\(\|p-q\|_{\rm TV}\le e^{\operatorname{osc}(f-f')}-1\).

\emph{Proof.}
\(p/q=e^{f-f'}Z_{f'}/Z_f\le e^{\sup(f-f')}e^{\sup(f'-f)}=e^{\operatorname{osc}(f-f')}\),
so \(p-q\le(e^{\operatorname{osc}}-1)q\) and
\(\|p-q\|_{\rm TV}=\int(p-q)_+\le e^{\operatorname{osc}}-1\).
\(\square\)

\textbf{Lemma 2.} If \(\omega,\omega'\) differ only at a link \(\ell'\)
sharing a plaquette with \(\ell\), then
\(\operatorname{osc}_U(f_\omega-f_{\omega'})\le4\beta_W\).

\emph{Proof.} One staple changes, \(M_\ell-M_\ell'=S_p-S_p'\) with
\(S_p,S_p'\in SU(N)\), so
\(|f_\omega(U)-f_{\omega'}(U)|=(\beta_W/N)|\operatorname{Re}\operatorname{tr}(U(S_p-S_p'))|\le(\beta_W/N)\|S_p-S_p'\|_1\le(\beta_W/N)\cdot2N\),
the trace norm of a difference of two unitaries being at most \(2N\);
the oscillation is at most twice the supremum. \(\square\)

If \(\ell'\) shares no plaquette with \(\ell\), the conditional
distribution at \(\ell\) does not depend on it. Each link lies in six
plaquettes with three further links each, and the six plaquettes share
no link besides \(\ell\), so there are exactly \(18\) links \(\ell'\)
with \(\rho_{\ell\ell'}\ne0\). Hence
\[\alpha=\sup_\ell\sum_{\ell'}\rho_{\ell\ell'}\le18\big(e^{4\beta_W}-1\big),\]
and \(\alpha<1\) for \(\beta_W<0.01352\), i.e.~\(g^2>2N/0.01352=148N\).

\hypertarget{from-the-condition-to-the-gap}{%
\subsection{3. From the condition to the
gap}\label{from-the-condition-to-the-gap}}

Dobrushin's theorem: \(\alpha<1\) implies a unique Gibbs state on
\(\mathbb Z^4\) and on every torus. The covariance estimate of Föllmer
(J. Funct. Anal. 46 (1982) 387; Künsch, loc. cit., for the decay
statements): under \(\alpha<1\), for local observables \(F,G\),
\[|\langle F;G\rangle|\ \le\ \sum_{\ell\in\operatorname{supp}F,\ \ell'\in\operatorname{supp}G}\delta_\ell F\,\big[(1-\rho)^{-1}\big]_{\ell\ell'}\,\delta_{\ell'}G ,\]
with \(\delta_\ell F\) the oscillation of \(F\) in the link \(\ell\),
and since \(\rho^n_{\ell\ell'}\) vanishes unless the link distance is at
most \(n\) lattice units,
\([(1-\rho)^{-1}]_{\ell\ell'}\le\alpha^{d(\ell,\ell')}/(1-\alpha)\).
Truncated correlations at time separation \(T\) therefore decay at least
like \(\alpha^{T/a}\), rate \(m=\log(1/\alpha)/a\). The Wilson transfer
matrix is positive by link reflection positivity, and for a local \(A\)
orthogonal to the vacuum
\(\langle A,\mathcal T^nA\rangle\le C_A\alpha^n\) forces its spectral
measure to vanish above \(\alpha\), so on the cyclic subspace of local
observables
\[\Delta_W=-\frac{\hbar c}{a}\log\|\mathcal T|_{\Omega^\perp}\|\ \ge\ \frac{\hbar c}{a}\log\frac1\alpha
\ \ge\ \frac{\hbar c}{a}\log\frac{1}{18(e^{24/g^2}-1)}\qquad(SU(3)).\]

\begin{longtable}[]{@{}
  >{\raggedright\arraybackslash}p{(\columnwidth - 4\tabcolsep) * \real{0.3333}}
  >{\raggedright\arraybackslash}p{(\columnwidth - 4\tabcolsep) * \real{0.3333}}
  >{\raggedright\arraybackslash}p{(\columnwidth - 4\tabcolsep) * \real{0.3333}}@{}}
\toprule\noalign{}
\begin{minipage}[b]{\linewidth}\raggedright
route
\end{minipage} & \begin{minipage}[b]{\linewidth}\raggedright
rigorous threshold (\(SU(3)\))
\end{minipage} & \begin{minipage}[b]{\linewidth}\raggedright
gap bound
\end{minipage} \\
\midrule\noalign{}
\endhead
\bottomrule\noalign{}
\endlastfoot
polymer expansion, all representations & \(g^2\ge1056\) &
\((\hbar c/a)\,4\log(g^2/1056)\) \\
Dobrushin single-site condition & \(g^2>444\) &
\((\hbar c/a)\,\log(g^2/432)\) \\
\end{longtable}

The two are complementary: the Dobrushin route enters earlier, the
polymer route's rate is larger once both apply. For \(U(1)\) the same
computation gives \(g^2>148\) with \(N=1\).

\hypertarget{the-block-criterion-as-a-specification}{%
\subsection{4. The block criterion, as a
specification}\label{the-block-criterion-as-a-specification}}

\begin{quote}
\textbf{Corrected.} The single-link form below is emptied by Elitzur's
theorem: the marginal of any link with an interior endpoint is Haar for
every boundary condition, so \(\rho_V(x,y)=0\) there and the sum sees
only the boundary layer. The form to verify is the mixing condition on
sub-blocks, organized around the conditional distribution of small
Wilson loops, and the boxes needed are of side \(3\) to \(5\) at
\(\beta_W\simeq5.7\); see \href{confinement-scale-bands.md}{the bands
note}.
\end{quote}

The Dobrushin--Shlosman constructive criterion replaces the single link
by a finite box \(V\) of links. Writing \(\mu_V(\cdot|\omega)\) for the
Gibbs measure in \(V\) with boundary condition \(\omega\) and, for
\(x\in V\) and \(y\notin V\),
\[\rho_V(x,y)=\sup_{\omega,\omega'\ \text{differ only at }y}\big\|\mu_V(\cdot|\omega)|_x-\mu_V(\cdot|\omega')|_x\big\|_{\rm TV},\]
the criterion is that for some finite \(V\)
\[\sum_{x\in V}\sum_{y\notin V}\rho_V(x,y)\ <\ |V| ,\] up to the
normalization convention of the original, and it implies the same
conclusions as the single-site condition, which is the case
\(V=\{\ell\}\). Its content is that the total influence of the boundary
on the box, summed over the box, is less than the size of the box; in a
mixing phase this holds once the side of \(V\) exceeds a few correlation
lengths, which is why it can hold at couplings where the single-site
condition fails. Dobrushin and Shlosman's complete analyticity
(\emph{Statistical Physics and Dynamical Systems}, Birkhäuser 1985, 371;
J. Stat. Phys. 46 (1987) 983) is the statement that some such \(V\)
exists, together with its equivalent forms.

\textbf{What a verification computes.} For a box \(V\) of side \(R\)
(about \(3R^4\) links) and a coupling \(\beta_W\): for each boundary
link \(y\) and each interior link \(x\), a rigorous upper bound on
\(\rho_V(x,y)\), which is a total-variation distance between two
marginals of Gibbs measures on \(SU(3)^{|V|}\); then the double sum
against \(|V|\). Lemma 1 bounds each \(\rho_V(x,y)\) by
\(e^{\operatorname{osc}}-1\) of the log-density difference of the two
marginals, but for \(|V|>1\) that oscillation involves the integral over
the other links of \(V\) and is what has to be bounded rigorously. The
openness radius in the interaction norm
\(\sup_\ell\sum_{p\ni\ell}\|\Phi_p\|_\infty\) follows from the margin
\(|V|-\sum\rho_V\) by the same lemma.

\textbf{Where it would have to hold.} The numerical evidence places the
\(SU(3)\) crossover at \(\beta_W\simeq5\)--\(6.5\) with \(\xi/a\) of
order \(1\) to \(10\); a verification at those couplings would need
\(R\) of order \(10\)--\(30\), that is \(10^4\)--\(10^6\) links, and
rigorous bounds on the corresponding integrals. That is far beyond
present practice, and this repository performs no numerics. The
specification is nonetheless complete: the quantity, the boxes, the
inequality, and the two lemmas that convert its margin into a gap and
into an openness radius.

\hypertarget{consequence-for-state}{%
\subsection{5. Consequence for STATE}\label{consequence-for-state}}

The strong-coupling side now has a third, elementary proof with the best
rigorous threshold for the Wilson transfer matrix, \(g^2>444\) for
\(SU(3)\), and the finite verification of the intermediate region has a
precise statement: the Dobrushin--Shlosman block criterion for the
Wilson interaction, whose single-link case is proved here. The two
numbers of the finite-verification note are unchanged, and the second of
them, \(g_{\rm DS}^2\), is exactly the smallest coupling at which the
block criterion has been verified for some \(V\).
\end{document}
